% paper/sections/11c_time_accuracy.tex
%
% §11.3  NS 方程式の総合精度
% ═══════════════════════════════════════════════════════════════════════════════
\subsection{NS 方程式の総合精度}
\label{sec:ns_time_accuracy}
% ═══════════════════════════════════════════════════════════════════════════════

§\ref{sec:ns_conservation} で保存則の物理的整合性を確認した．
本節では一歩進めて，NS ソルバ全体の\textbf{収束次数}——
時間精度 $\Ord{\Delta t^p}$ と空間精度 $\Ord{h^q}$ ——を定量的に確認する．
コンポーネント単体（第~\ref{sec:numerical_integrity}章）の精度が
パイプライン全体でどの程度発現するかが焦点である．

% ─────────────────────────────────────────────────────────────────────────────
\subsubsection{Taylor--Green 渦による時間収束次数}
\label{sec:tgv_temporal}
% ─────────────────────────────────────────────────────────────────────────────

\noindent\textbf{目的：}
AB2+IPC スキームが NS ソルバ全体で設計時間精度 $\Ord{\Delta t^2}$ を達成することを確認する．

\noindent\textbf{評価対象：}
Predictor（AB2 + CN）+ PPE + Corrector の結合時間精度．

\noindent\textbf{手順とパラメタ：}
\begin{itemize}[noitemsep]
  \item §\ref{sec:energy_conservation} と同一の Taylor--Green 渦（$\mathit{Re}=100$）
  \item $N = 64$ に固定（空間誤差を十分に小さくする）
  \item $\Delta t$ を系統的に変化させ（$n = 5, 10, 20, 40, 80, 160$ ステップ），
        $T = 0.5$ における $L^\infty(\bu)$ を計測
\end{itemize}

\noindent\textbf{結果：}

\begin{table}[htbp]
\centering
\caption{NS 時間精度：Taylor--Green 渦（$N=64$，$\mathit{Re}=100$，AB2+IPC）}
\label{tab:tgv_temporal}
\renewcommand{\arraystretch}{1.3}
\begin{tabular}{rrrc}
\toprule
ステップ数 & $\Delta t$ & $L^\infty(\bu)$ & 次数 \\
\midrule
  5  & $0.1000$ & $3.35 \times 10^{-9}$  & --- \\
 10  & $0.0500$ & $8.37 \times 10^{-10}$ & $2.00$ \\
 20  & $0.0250$ & $2.10 \times 10^{-10}$ & $1.99$ \\
 40  & $0.0125$ & $5.33 \times 10^{-11}$ & $1.98$ \\
 80  & $0.0063$ & $1.42 \times 10^{-11}$ & $1.91$ \\
160  & $0.0031$ & $4.34 \times 10^{-12}$ & $1.71$ \\
\bottomrule
\end{tabular}
\end{table}

\begin{figure}[htbp]
\centering
\includegraphics[width=0.62\linewidth]{figures/ch11_tgv_temporal}
\caption{NS 時間精度：Taylor--Green 渦の $L^\infty$ 速度誤差 vs $\Delta t$（$N=64$，$\mathit{Re}=100$）．
$\Delta t \geq 0.0125$ で理論通りの $\Ord{\Delta t^2}$ 勾配に乗る．
$\Delta t \leq 0.006$ では丸め誤差のフロアにより見かけ上の次数が低下．}
\label{fig:tgv_temporal}
\end{figure}

\noindent\textbf{結論：}
$\Delta t \geq 0.0125$ で精確な $\Ord{\Delta t^2}$ 収束を確認した
（次数 $1.98$--$2.00$；図~\ref{fig:tgv_temporal}）．
AB2（$\Ord{\Delta t^2}$）+ IPC 分割誤差（$\Ord{\Delta t^2}$，van Kan 1986）の理論と一致する．
$\Delta t \leq 0.006$ では丸め誤差のフロア（$\sim 10^{-12}$）に接近し見かけ上の次数が低下する．

% ─────────────────────────────────────────────────────────────────────────────
\subsubsection{Kovasznay 流による空間精度の検証}
\label{sec:kovasznay}
% ─────────────────────────────────────────────────────────────────────────────

\noindent\textbf{目的：}
CCD 離散化した NS 方程式の空間収束次数を，
非線形対流項を含む定常解析解を用いて検証する．

\noindent\textbf{評価対象：}
CCD 離散化 NS 残差の格子収束特性（非線形項 + 境界スキームの影響を含む全体精度）．

\noindent\textbf{手順とパラメタ：}
\begin{itemize}[noitemsep]
  \item Kovasznay (1948) の 2 次元定常 NS 解析解（$\mathit{Re} = 40$，
        $\lambda = \mathit{Re}/2 - \sqrt{\mathit{Re}^2/4 + 4\pi^2} \approx -6.42$）：
        \begin{align}
          u &= 1 - e^{\lambda x}\cos(2\pi y), &
          v &= \frac{\lambda}{2\pi}\,e^{\lambda x}\sin(2\pi y), &
          p &= -\tfrac{1}{2}\,e^{2\lambda x}
        \end{align}
  \item 計算領域 $[0, 1] \times [-0.5, 0.5]$，ディリクレ境界条件
  \item 格子 $N = 16, 32, 64, 128$
  \item 評価方法：CCD 離散化 NS 残差 $R = -\bu \cdot\bnabla\bu + \nu\bnabla^2\bu - \bnabla p$
        を解析解上で評価（MMS 残差法）
\end{itemize}

\noindent\textbf{結果：}

\begin{table}[htbp]
\centering
\caption{Kovasznay 流の CCD 離散化 NS 残差（$\mathit{Re}=40$，ディリクレ BC）}
\label{tab:kovasznay}
\renewcommand{\arraystretch}{1.3}
\begin{tabular}{rrcrc}
\toprule
$N$ & NS 残差 $\|R\|_\infty$ & 次数 & $\|\bnabla\cdot\bu\|_\infty$ & 次数 \\
\midrule
 16 & $2.60 \times 10^{-4}$ & --- & $6.60 \times 10^{-5}$ & --- \\
 32 & $2.28 \times 10^{-5}$ & $3.51$ & $1.24 \times 10^{-6}$ & $5.74$ \\
 64 & $1.55 \times 10^{-6}$ & $3.88$ & $2.05 \times 10^{-8}$ & $5.91$ \\
128 & $9.93 \times 10^{-8}$ & $3.96$ & $3.27 \times 10^{-10}$ & $5.97$ \\
\bottomrule
\end{tabular}
\end{table}

\begin{figure}[htbp]
\centering
\includegraphics[width=0.62\linewidth]{figures/ch11_kovasznay_convergence}
\caption{Kovasznay 流の CCD 離散化 NS 残差の格子収束（$\mathit{Re}=40$）．
NS 運動量残差（青丸）は $\Ord{h^{3.8}}$，
連続方程式残差（赤四角）は $\Ord{h^{5.9}}$ で収束する．}
\label{fig:kovasznay_convergence}
\end{figure}

\noindent\textbf{結論：}
連続方程式残差は $\Ord{h^{5.97}}$（CCD 内部精度 $\Ord{h^6}$ に整合；図~\ref{fig:kovasznay_convergence}）．
NS 運動量残差は $\Ord{h^{3.96}}$：非線形対流項 $\bu \cdot \bnabla\bu$ がディリクレ境界近傍で
CCD 境界スキーム（$\Ord{h^5}$）の誤差を非線形に増幅するため，
実効精度は $\Ord{h^4}$ 程度となる．
$N=128$ で NS 残差は $10^{-8}$ に達し，十分な空間精度を確認した．
